% ============================================================================
%  04-07: INTEGRACIÓN CON LA APLICACIÓN ASP.NET CORE (.NET 10)
% ============================================================================

\section{Integraci\'on con el backend ASP.NET Core (.NET 10)}

\subsection{Mapeo multiesquema y restricciones en \texttt{ApplicationDbContext}}

Para sincronizar al 100\,\% el modelo de dominio C\# con la base de datos desplegada en Supabase, se refactoriz\'o \texttt{ApplicationDbContext} utilizando Fluent API de Entity Framework Core 10 y el proveedor \texttt{Npgsql}.
Las tablas de ASP.NET Core Identity permanecen en \texttt{public} (enlazando \texttt{ApplicationUser} con \texttt{Sucursal} y \texttt{Region}), mientras que las catorce entidades de negocio y auditor\'ia se mapean expl\'icitamente a los esquemas \texttt{catalogo}, \texttt{ventas}, \texttt{operaciones} y \texttt{auditoria}, incluyendo la entidad \texttt{MermaInventario} para el control de obsolescencia (ODS~12).

\begin{lstlisting}[language={[Sharp]C}, caption={Extracto real de ApplicationDbContext.cs con mapeo multiesquema y restricciones}, label=lst:efcore]
public class ApplicationDbContext(DbContextOptions<ApplicationDbContext> options)
    : IdentityDbContext<ApplicationUser>(options)
{
    // -- Esquema catalogo (3FN) --
    public DbSet<Categoria> Categorias { get; set; }
    public DbSet<Marca> Marcas { get; set; }
    public DbSet<Componente> Componentes { get; set; }
    public DbSet<Combo> Combos { get; set; }
    public DbSet<ComboComponente> ComboComponentes { get; set; }

    // -- Esquema ventas --
    public DbSet<Pedido> Pedidos { get; set; }
    public DbSet<PedidoItem> PedidoItems { get; set; }
    public DbSet<CarritoAbandonado> CarritosAbandonados { get; set; }

    // -- Esquema operaciones --
    public DbSet<Region> Regiones { get; set; }
    public DbSet<Sucursal> Sucursales { get; set; }
    public DbSet<Proveedor> Proveedores { get; set; }

    // -- Esquema auditoria (incluye ODS 12) --
    public DbSet<AuditoriaComponente> AuditoriaComponentes { get; set; }
    public DbSet<AuditoriaPedido> AuditoriaPedidos { get; set; }
    public DbSet<MermaInventario> MermasInventario { get; set; }

    protected override void OnModelCreating(ModelBuilder modelBuilder)
    {
        base.OnModelCreating(modelBuilder);

        modelBuilder.Entity<Componente>(e =>
        {
            e.ToTable("Componentes", "catalogo", t =>
            {
                t.HasCheckConstraint("ck_componentes_precio", "\"Precio\" > 0");
                t.HasCheckConstraint("ck_componentes_stock", "\"Stock\" >= 0");
                t.HasCheckConstraint("ck_componentes_min", "\"StockMinimo\" >= 0");
                t.HasCheckConstraint("ck_componentes_obsol",
                    "\"EstadoObsolescencia\" IN ('Vigente', 'PorDescontinuar', 'Obsoleto', 'Merma')");
            });
            e.HasIndex(c => c.CategoriaId).HasDatabaseName("idx_comp_categoria");
            e.HasIndex(c => c.MarcaId).HasDatabaseName("idx_comp_marca");
            e.HasIndex(c => c.Precio).HasDatabaseName("idx_comp_precio");
        });

        modelBuilder.Entity<PedidoItem>(e =>
        {
            e.ToTable("PedidoItems", "ventas", t =>
            {
                t.HasCheckConstraint("ck_pi_exclusivo",
                    "(\"ComponenteId\" IS NOT NULL AND \"ComboId\" IS NULL) OR " +
                    "(\"ComboId\" IS NOT NULL AND \"ComponenteId\" IS NULL)");
            });
        });

        modelBuilder.Entity<CarritoAbandonado>(e =>
        {
            e.ToTable("CarritosAbandonados", "ventas");
            e.Property(c => c.ItemsJson).HasColumnType("jsonb");
            e.HasIndex(c => c.ItemsJson).HasMethod("gin").HasDatabaseName("idx_carritos_json");
        });

        modelBuilder.Entity<MermaInventario>(e =>
        {
            e.ToTable("MermasInventario", "auditoria", t =>
            {
                t.HasCheckConstraint("ck_merma_cant", "\"Cantidad\" > 0");
                t.HasCheckConstraint("ck_merma_motivo",
                    "\"Motivo\" IN ('Obsolescencia', 'DanoFisico', 'Garantia', 'ReciclajeRAEE', 'Otro')");
                t.HasCheckConstraint("ck_merma_disp",
                    "\"DisposicionFinal\" IN ('ReciclajeCertificado', 'Reacondicionado', 'Donacion', 'BajaDefinitiva')");
            });
        });
    }
}
\end{lstlisting}

\subsection{Invocaci\'on de \texttt{ventas.sp\_confirmar\_venta} en \texttt{TiendaController}}

En lugar de insertar pedidos y descontar stock desde memoria sin bloqueo de filas, el m\'etodo \texttt{ProcesarPago} de \texttt{TiendaController} serializa el carrito a \texttt{JSONB} e invoca directamente \texttt{CALL ventas.sp\_confirmar\_venta} mediante \texttt{NpgsqlCommand}.
Si PostgreSQL detecta falta de stock y ejecuta \texttt{ROLLBACK}, la excepci\'on \texttt{PostgresException} es capturada por el controlador y presentada de forma clara al cliente en la vista de \textit{checkout}.

\begin{lstlisting}[language={[Sharp]C}, caption={Invocacion real de CALL ventas.sp\_confirmar\_venta en TiendaController.cs}, label=lst:tienda-sp]
[HttpPost, Authorize]
[ValidateAntiForgeryToken]
public async Task<IActionResult> ProcesarPago(MetodoPago metodoPago, string? notas)
{
    var carrito = ObtenerCarrito();
    if (!carrito.Items.Any()) return RedirectToAction(nameof(Carrito));

    var user = await userManager.GetUserAsync(User);
    if (user == null) return Challenge();

    var fechaHoy   = DateTime.UtcNow;
    var numPedido  = $"TPC-{fechaHoy:yyyyMMdd-HHmmss}-{Random.Shared.Next(100, 999)}";
    var sucursalId = user.SucursalId ??
        await db.Sucursales.Where(s => s.Activa).Select(s => (int?)s.Id).FirstOrDefaultAsync();

    var itemsJson = JsonSerializer.Serialize(carrito.Items.Select(i => new
    {
        tipo = i.Tipo,
        id = i.Id,
        cantidad = i.Cantidad
    }));

    try
    {
        var conn = (NpgsqlConnection)db.Database.GetDbConnection();
        bool debeCerrar = conn.State != ConnectionState.Open;
        if (debeCerrar) await conn.OpenAsync();

        int pedidoId;
        try
        {
            await using var cmd = new NpgsqlCommand(
                "CALL ventas.sp_confirmar_venta(@p_user_id, @p_numero_pedido, @p_items, " +
                "@p_metodo_pago, @p_sucursal_id, @p_notas, @p_pedido_id)", conn);

            cmd.Parameters.Add(new NpgsqlParameter("p_user_id", NpgsqlDbType.Text) { Value = user.Id });
            cmd.Parameters.Add(new NpgsqlParameter("p_numero_pedido", NpgsqlDbType.Varchar) { Value = numPedido });
            cmd.Parameters.Add(new NpgsqlParameter("p_items", NpgsqlDbType.Jsonb) { Value = itemsJson });
            cmd.Parameters.Add(new NpgsqlParameter("p_metodo_pago", NpgsqlDbType.Smallint) { Value = (short)metodoPago });
            cmd.Parameters.Add(new NpgsqlParameter("p_sucursal_id", NpgsqlDbType.Integer) { Value = (object?)sucursalId ?? DBNull.Value });
            cmd.Parameters.Add(new NpgsqlParameter("p_notas", NpgsqlDbType.Text) { Value = (object?)notas ?? DBNull.Value });

            var outPedidoId = new NpgsqlParameter("p_pedido_id", NpgsqlDbType.Integer)
            {
                Direction = ParameterDirection.InputOutput,
                Value = DBNull.Value
            };
            cmd.Parameters.Add(outPedidoId);

            await cmd.ExecuteNonQueryAsync();
            pedidoId = Convert.ToInt32(outPedidoId.Value);
        }
        finally
        {
            if (debeCerrar) await conn.CloseAsync();
        }

        HttpContext.Session.Remove(CarritoKey);
        TempData["Exito"] = $"Compra completada con integridad transaccional. Pedido {numPedido}";
        return RedirectToAction(nameof(Factura), new { id = pedidoId });
    }
    catch (PostgresException pgEx)
    {
        TempData["Error"] = $"No se pudo confirmar la compra (ROLLBACK): {pgEx.MessageText}";
        return RedirectToAction(nameof(Checkout));
    }
}
\end{lstlisting}
